\chapter{Intertwiner чётного coherence-угла и семейного триплета}
\label{ch:v8-baryon-c0-coherence-even-corner-family-triplet-intertwiner}

\section{Текущие представления}

Алгебраическая форма портала использовала совпадение размерностей, но
физический morphism обязан сплетать все объявленные действия. Чётный угол
coherence-chain имеет тип
\begin{equation}
 E_{\rm coh}=\Lambda^2W^*,
 \qquad
 W=W_{eX}\oplus W_Y\simeq\mathbb C^2\oplus\mathbb C,
 \qquad
 E_{\rm coh}\simeq\mathbf1\oplus\mathbf2
 \label{eq:v8-coherence-family-intertwiner-source-type}
\end{equation}
относительно $G_{\rm ch}=U(2)_{eX}\times U(1)_Y$. Семейный target равен
\begin{equation}
 E_{\rm fam}\simeq\mathbf3_{SO(3)_{\rm fam}}.
 \label{eq:v8-coherence-family-intertwiner-target-type}
\end{equation}
В текущем контракте $SO(3)_{\rm fam}$ действует тривиально на
$E_{\rm coh}$, а $G_{\rm ch}$ --- тривиально на $E_{\rm fam}$. Поэтому
полная группа имеет вид
\begin{equation}
 G_{\rm current}=G_{\rm ch}\times SO(3)_{\rm fam}.
 \label{eq:v8-coherence-family-intertwiner-product-group}
\end{equation}

\section{Нулевой Hom в текущей типизации}

Пусть $L_1,L_2,L_3$ --- стандартные вещественные генераторы $so(3)$, а
$Z:E_{\rm coh}\to E_{\rm fam}$. Семейные уравнения сплетения равны
\begin{equation}
 L_iZ-Z0=L_iZ=0,
 \qquad i=1,2,3.
 \label{eq:v8-coherence-family-current-so3-system}
\end{equation}
На девяти вещественных неизвестных их матрица имеет
\begin{equation}
 \rank\mathcal C_{SO(3)}=9,
 \qquad
 \dim\ker\mathcal C_{SO(3)}=0,
 \qquad
 \operatorname{Hom}_{SO(3)}(\mathbf1^{\oplus3},\mathbf3)=0.
 \label{eq:v8-coherence-family-current-so3-hom}
\end{equation}
Ограничение до тетраэдрического остатка не помогает. Для циклической
перестановки $R_3$ и полуоборота $R_2=\operatorname{diag}(1,-1,-1)$
получается
\begin{equation}
 (R_3-I)Z=(R_2-I)Z=0,
 \qquad
 \rank\mathcal C_{A_4}=9,
 \qquad
 \operatorname{Hom}_{A_4}(\mathbf1^{\oplus3},\mathbf3)=0.
 \label{eq:v8-coherence-family-current-a4-hom}
\end{equation}

Независимый канальный тест также даёт нуль. Если
$\rho_{\wedge^2}$ обозначает действие $G_{\rm ch}$ на
$\Lambda^2W^*$, то
\begin{equation}
 0Z-Z\rho_{\wedge^2}(X)=0
 \quad\forall X\in\mathfrak u(2)\oplus\mathfrak u(1)
 \quad\Longrightarrow\quad Z=0.
 \label{eq:v8-coherence-family-current-channel-hom}
\end{equation}
Следовательно,
\begin{equation}
 \operatorname{Hom}_{G_{\rm current}}(E_{\rm coh},E_{\rm fam})=0.
 \label{eq:v8-coherence-family-current-product-hom}
\end{equation}
Это типизационный no-go, а не отсутствие удачного выбора базиса.

\section{Единственная условная Hodge-лазейка}

Если условно заменить канальное разложение единым ориентированным
вещественным триплетом
\begin{equation}
 W_{\rm ext}\simeq\mathbb R^3_{SO(3)_{\rm diag}},
 \qquad
 E_{\rm coh}^{\rm ext}=\Lambda^2\mathbb R^3,
 \qquad
 E_{\rm fam}=\mathbb R^3,
 \label{eq:v8-coherence-family-promoted-types}
\end{equation}
то метрика и ориентация задают Hodge-map
\begin{equation}
 *=\begin{pmatrix}
 0&0&1\\
 0&-1&0\\
 1&0&0
 \end{pmatrix}:
 \Lambda^2\mathbb R^3\longrightarrow\mathbb R^3.
 \label{eq:v8-coherence-family-hodge-map}
\end{equation}
В базисе $(e_1\wedge e_2,e_1\wedge e_3,e_2\wedge e_3)$ он удовлетворяет
\begin{equation}
 L_i*=*\,\rho_{\wedge^2}(L_i),
 \qquad i=1,2,3.
 \label{eq:v8-coherence-family-hodge-intertwining}
\end{equation}
Полная система имеет ранг восемь:
\begin{equation}
 \rank\mathcal C_{\rm diag}=8,
 \qquad
 \operatorname{Hom}_{SO(3)_{\rm diag}}
 (\Lambda^2\mathbb R^3,\mathbb R^3)=\mathbb R*.
 \label{eq:v8-coherence-family-promoted-hom}
\end{equation}
Нормировка метрики оставляет только два знака,
\begin{equation}
 *^{\mathsf T}*=I_3,
 \qquad
 \det *=1,
 \qquad
 (c*)^{\mathsf T}(c*)=I_3
 \Longleftrightarrow c=\pm1.
 \label{eq:v8-coherence-family-hodge-normalization}
\end{equation}
Ориентация выбирает знак лишь после того, как сама происходит из
родительской структуры.

\section{Цена условного повышения}

Текущее $W$ не является тремя копиями одного бимодуля: $W_Y$ имеет тип,
отличный от $W_e,W_X$. Поэтому непрерывное смешивание всех трёх координат
требует
\begin{equation}
 W_e\simeq W_X\simeq W_Y
 \quad\text{как бимодули},
 \qquad
 SO(3)_{\rm diag}\subset U(3)_{\rm accidental},
 \label{eq:v8-coherence-family-bimodule-promotion-cost}
\end{equation}
а также диагонального отождествления этой новой группы с
$SO(3)_{\rm fam}$. Классификация конечных спектральных троек требует
сохранять именно бимодульные вершины и рёбра, а не только размерности
\cite{v8-coherence-family-krajewski}; поэтому такое повышение является
новой архитектурой.

Реестры имеют вид
\begin{equation}
 \begin{aligned}
 N_{\rm current\ intertwiner}&=0/4,
 &N_{\rm conditional\ shape}&=4/4,\\
 N_{\rm new\ origin}&=0/3,\\
 \mathcal M_{\rm new}&=\text{bimodule, diagonal lock, orientation}.
 \end{aligned}
 \label{eq:v8-coherence-family-intertwiner-ledgers}
\end{equation}

\begin{theorem}[Текущий Hom равен нулю]
Для действующей product-group
$\bigl(U(2)_{eX}\times U(1)_Y\bigr)\times SO(3)_{\rm fam}$ не существует
ненулевого morphism из coherence-угла $\Lambda^2W^*$ в стандартный
семейный триплет. Запрет сохраняется после ограничения семейной группы до
$A_4$.
\end{theorem}

\begin{proposition}[Условная единственность Hodge-map]
После повышения $W$ до ориентированного стандартного триплета диагонального
$SO(3)$ пространство intertwiner равно $\mathbb R*$, а изометрические карты
равны $\pm *$.
\end{proposition}

\begin{nogocriterion}[Запрет случайного $U(3)$]
Нельзя объявлять условный Hodge-map существующим физическим morphism, пока
различные канальные бимодули не объединены одним родителем и не выведено
диагональное family--channel действие. Равенство размерностей и случайная
$U(3)$-симметрия нормы этого не доказывают.
\end{nogocriterion}

\begin{protocol}\begin{enumerate}
\item текущая группа: \textbf{$G_{\rm ch}\times SO(3)_{\rm fam}$};
\item текущий $SO(3)$-Hom: \textbf{$0$};
\item текущий $A_4$-Hom: \textbf{$0$};
\item текущий канальный Hom: \textbf{$0$};
\item полный текущий intertwiner: \textbf{$0/4$};
\item условный диагональный Hom: \textbf{$\mathbb R*$};
\item изометрические карты: \textbf{$\pm *$};
\item условная форма: \textbf{$4/4$};
\item новых parent-структур: \textbf{$3$};
\item portal parent-origin: \textbf{не закрыт};
\item следующий гейт: \textbf{бимодульная совместимость channel-triplet повышения}.
\end{enumerate}\end{protocol}

Машинный сертификат сохранён в
\path{s2t_v8_baryon_c0_singlet_triplet_central_gap_coherence_even_corner_family_triplet_intertwiner_gate_results.json}.

\begin{thebibliography}{9}
\bibitem{v8-coherence-family-krajewski}
T.~Krajewski,
\emph{Classification of Finite Spectral Triples},
arXiv:hep-th/9701081.
\bibitem{v8-coherence-family-quiver}
M.~Harada, G.~Wilkin,
\emph{Morse theory of the moment map for representations of quivers},
arXiv:0807.4734.
\end{thebibliography}